% !TeX root = ./chapter_1_main.tex
\section{Optical Dipole Trap}\label{sec:ODT_theory}
Optical traps are experimental tools that allow the confinement of atom or ions within a localized region of space, under the condition that these particles have a kinetic energy lower than the trap potential depth (typically below $1$mK). For the work presented in this thesis, among all possible trapping techniques, the focus will be strictly placed on Optical Dipole Traps (ODTs) for neutral atoms.

Optical dipole traps rely on the electric dipole interaction of the atoms with far-detuned light. Following the work of \citeauthor*{grimm} \cite{grimm}, this section reviews the theory behind the working mechanism of this technique, particularly in the case of red-detuned ODTs. 

%* ##########################################################
\subsection{Oscillator model}
%* ##########################################################
The optical dipole force arises from the interaction of the induced electric dipole moment with the intensity gradient of the light field. It is a conservative force, thus in the following its potential energy will be derived modelling the atom as an oscillator subjected to a classical oscillating electric field.

\subsubsection{Interaction between induced dipole moment and driving field}

%! Magari qua ci sta eslicitare le assunzioni

The laser field is modeled as the oscillating electric field: 
\begin{equation}
    \mathbf{E}(\mathbf{r},t) = \hat{\mathbf{e}} \mathcal{E}(\mathbf{r}) e^{-i\omega t} + c.c. \,\, ,
\end{equation}

where $\mathbf{\hat{e}}$ is the unitary polarization vector and $\mathcal{E}(r)$ is the complex amplitude of the field. Denoting the atomic polarizability as $\alpha$, the induced dipole momentum of the atom is given by:
\begin{equation}
    \mathbf{p}(\mathbf{r}, t) = \alpha \mathbf{E}(\mathbf{r},t) = \hat{\mathbf{e}} \alpha \mathcal{E}(\mathbf{r})e^{-i\omega t} + c.c. 
\end{equation}
In the semiclassical framework, the interaction of a neutral atom with a laser light field can be described using the dipole approximation. The interaction Hamiltonian takes the form $H_\text{dip} = - \mathbf{p \cdot E}$. For an induced dipole, where the dipole moment depends linearly on the external field, the potential energy is obtained by integrating the work required to induce the dipole as the electric field is increased from zero to its final value:
\begin{equation}
    U_\text{dip}(t) = \int_0^\mathbf{E} -\alpha \mathbf{E'}\cdot d \mathbf{E'} = -\frac{1}{2} \alpha E^2 = - \frac{1}{2} \mathbf{p \cdot E} \,. 
\end{equation}

Averaging over a field period \cite{grimm}:
\begin{equation}
    \label{eq:dipole_potential_alpha}
    U_\text{dip}= - \frac{1}{2} \langle \mathbf{p \cdot E}\rangle= -\frac{1}{2\epsilon_0 c} \text{Re}[\alpha]I,
\end{equation}

where $I=2 \epsilon_0 c |\mathcal{E}|^2$ is the time-averaged field intensity.\\
This result shows that an atom in an oscillating laser field experiences a force proportional to the field intensity $I$, and to the real part of the polarizability, which represent the in-phase component of the dipole oscillation, responsible for the dispersive properties of the interaction.

It follows that the dipole force is a conservative force proportional to the electric field intensity gradient:
\begin{equation}
    \mathbf{F}_\text{dip} (\mathbf{r}) = -\nabla U_\text{dip} = \frac{1}{2\epsilon_0 c} \text{Re}[\alpha] \nabla I(\mathbf{r}).
\end{equation}
The dipole oscillation absorbs a power
\begin{equation}
    P_\text{abs} = \langle\mathbf{\dot{p}E}\rangle = 2\omega \text{Im}[\alpha \mathcal{E}^2] = \frac{\omega}{\epsilon_0 c}\text{Im}[\alpha] I,
\end{equation}
which is proportional to the out-of-phase component of the oscillating dipole. The absorbed power can be seen as series of absorption and spontaneous emission of photons with energy $\hbar\omega$. Therefore, the total scattering rate is given by
\begin{equation}
    \label{eq:scattering_rate_alpah}
    \Gamma_\text{sc} (\mathbf{r})=\frac{P_\text{abs}}{\hbar\omega}=\frac{1}{\hbar\epsilon_0 c}\text{Im} [\alpha] I(\mathbf{r}).
\end{equation}

\subsubsection{Atomic polarizability}
The polarizability $\alpha$ can be calculated by considering the classical Lorentz oscillator model for the atom, according to which the polarizability can be calculated by solving the differential equation
\begin{equation}
    \label{eq:Lorentz_polarizability}
    \ddot{x}+\Gamma(\omega) \dot{x} + \omega_0^2x = -\frac{eE(t)}{m_e},
\end{equation}
Using $E(t)=\mathcal{E}_0\exp(-i\omega t)$, thus looking for a solution in the form $x(t)=x_0\exp(-i\omega t)$, equation \eqref{eq:Lorentz_polarizability} gives
\begin{equation}
    \label{eq:Lorentz_polarizability_solution}
    x_0 = - \frac{e\mathcal{E}_0}{m_e} \,\frac{1}{\omega_0^2 -\omega^2 -i\omega\Gamma(\omega)}.
\end{equation}
Using $p= -ex(t)=-ex_0\exp(-i\omega t)$ and \eqref{eq:Lorentz_polarizability_solution}, one can easily obtain the polarizability
\begin{equation}
    \alpha(\omega) = \frac{e^2}{m_e} \,\frac{1}{\omega_0^2 -\omega^2 -i\omega\Gamma(\omega)},
\end{equation}
where $\Gamma(\omega)$ is the classical damping rate due to the radiative energy loss:
\begin{equation}
    \Gamma^\text{class.}(\omega) = \frac{e^2 \omega^2}{6 \pi \epsilon_0 m_e c^3}.
\end{equation}
Defining the on-resonance spontaneous emission rate $\Gamma\equiv  \Gamma^\text{class.}(\omega_0) = (\omega_0/\omega)^2\,\Gamma^\text{class.}(\omega)$, the polarizability becomes
\begin{equation}
    \label{eq:polarizability}
    \alpha(\omega) = 6 \pi \epsilon_0 c^3 \frac{\Gamma / \omega_0^2}{\omega_0^2 - \omega^2 - i (\omega^3 / \omega_0^2) \Gamma}.
\end{equation}

When considering a full quantum picture of a two level system subjected to classical radiation, $\Gamma$ should be substituted with the spontaneous emission rate
\begin{equation}
    \label{eq:quantum_gamma}
    \Gamma = \frac{\omega_0^3}{3\pi\epsilon_0\hbar c^3}|\langle e|\,\hat{d}\,|g\rangle|^2,
\end{equation}
where $\hat{d}$ is the dipole moment operator.\\
Nevertheless, optical dipole traps work in far-detuned low saturation regime. Consequently, the transitions to the excited state are strongly suppressed, and the atom can be considered as occupying the ground state only. In this limit, the atomic equation of motion reduces to a classical damped harmonic oscillator model described by equation \ref{eq:Lorentz_polarizability} (for a rigorous derivation see Chapter V of \citeauthor*{steck_quantum_optics} \cite{steck_quantum_optics}). Therefore, the classical polarizability given by equation \ref{eq:polarizability} represent a well suited approximation for ODTs.

\subsubsection{Dipole potential and scattering rate}
Substituting $\alpha$ from \eqref{eq:polarizability} in \eqref{eq:dipole_potential_alpha} and \eqref{eq:scattering_rate_alpah} one gets the expressions of the dipole potential and the scattering rate of a dipole trap as explicit functions of $\omega$, $\omega_0$ and $\Gamma$:
\begin{equation}
    U_{\text{dip}}(\mathbf{r}) = -\frac{3 \pi c^2}{2 \omega_0^3} \left( \frac{\Gamma}{\omega_0 - \omega} + \frac{\Gamma}{\omega_0 + \omega} \right) I(\mathbf{r}),
\end{equation}
\begin{equation}
    \Gamma_{\text{sc}}(\mathbf{r}) = \frac{3 \pi c^2}{2 \hbar \omega_0^3} \left( \frac{\omega}{\omega_0} \right)^3 \left( \frac{\Gamma}{\omega_0 - \omega} + \frac{\Gamma}{\omega_0 + \omega} \right)^2 I(\mathbf{r}).
\end{equation}

% These equations show two resonance frequencies for $\omega=\pm\omega_0$. Introducing the detuning $\Delta = \omega-\omega_0$, for the far-red regime it holds that $\Gamma\ll|\Delta|$. For instance, the $D_2$ line in Rb atoms has a transition frequency of $\sim385$~THz thus a trap typically works at few THz of detuning, while $\Gamma/(2\pi)\sim 6$~MHz. This allows the introduction of the well-known rotating-wave approximation, according to which the anti-resonant term can be neglected, further simplifying the equations to
These equations show two resonance frequencies for $\omega=\pm\omega_0$. Introducing the detuning $\Delta = \omega-\omega_0$, in typical experimental conditions, it holds that $|\Delta|\ll\omega_0$. This allows the introduction of the well-known rotating-wave approximation, according to which the anti-resonant term can be neglected, further simplifying the equations to 

\begin{equation}
    \label{eq:dipole_potential}
    U_{\text{dip}}(\mathbf{r}) = \frac{3 \pi c^2}{2 \omega_0^3} \frac{\Gamma}{\Delta} I(\mathbf{r})
\end{equation}
\begin{equation}
    \label{eq:dipole_scattering}
    \Gamma_{\text{sc}}(\mathbf{r}) = \frac{3 \pi c^2}{2 \hbar \omega_0^3} \left( \frac{\Gamma}{\Delta} \right)^2 I(\mathbf{r}) = \frac{\Gamma}{\hbar\Delta}U_\text{dip}
\end{equation}

Equations \eqref{eq:dipole_potential} and \eqref{eq:dipole_scattering} are the fundamental equations governing the behavior of optical dipole traps and yield important physical considerations for their implementation. The sign of $U_\text{dip}$ is determined by the sign of the detuning $\Delta$. For red detuning ($\Delta<0$) the potential is negative, confining atoms at intensity maxima, whereas for blue detuning ($\Delta>0$), the potential is positive pushing atoms to intensity minima. Because the potential depth scales as $I/\Delta$ while the scattering rate scales as $I/\Delta^2$, optical traps are typically realized with high intensity field and large detunings, in this way, the trap depth is maximized while maintaining a low scattering rate.



%* ##########################################################
\subsection{Multi-level atoms}
%* ##########################################################

Moving from the two-level atom to a real atom, the electronic transitions become more complex, and the calculation of the dipole potential must take into consideration all the possible allowed transitions. This lead to a dipole potential that depends on the particular electronic substate.

\subsubsection{Ground-state light shifts}
Labelling as $\varepsilon_i$ the energy levels of the multi-level atom, the effect of the laser light on the levels can be described by means of the non-degenerate time-independent second order perturbation theory. Considering the electric dipole interaction Hamiltonian $\mathcal{H}_I=-\hat{\mathbf{d}}\cdot\mathbf{E}$, the perturbative correction on the $i$-th energy level is
\begin{equation}
    \label{eq:energy_perturbative_correction}
    \Delta E_i = \sum_{j \neq i} \frac{|\langle j | \mathcal{H}_I | i\rangle|^2}{\varepsilon_i - \varepsilon_j}.
\end{equation}

The calculation can be developed by employing the dressed atom approach, in which the quantum states describe the combined atom-field system. Defining the unperturbed atomic ground state energy as zero, in this framework an unperturbed dressed state has an energy $\varepsilon_i=n\hbar\omega$ that depends solely on the number of photons $n$ in the field. The absorption of a photon transitions the system to another dressed state characterized by the energy $\varepsilon_j=\hbar\omega_{j} + (n-1)\hbar\omega$, where $\hbar\omega_{j}$ represent the bare atomic transition energy. In such a way, the denominator of equation \eqref{eq:energy_perturbative_correction} reduces to
\begin{equation}
    \label{eq:perturbative_expantion_denominator}
    \varepsilon_i - \varepsilon_j = -\hbar\omega_j + \hbar\omega = -\hbar\Delta_{ij}, 
\end{equation}
where $\Delta_{ij}$ is the detuning of the $i \rightarrow j$ transition with respect to the laser field.

Considering a two-level atom first, \eqref{eq:energy_perturbative_correction} has only one term, using equations \eqref{eq:quantum_gamma} and \eqref{eq:perturbative_expantion_denominator} and $I=2\epsilon_0c|\mathcal{E}|^2$, one gets:
\begin{equation}
    \label{eq:light_shift_2lv_atoms}
    \Delta E = \pm \frac{|\langle e | \mu | g \rangle|^2}{\Delta} |E|^2 = \pm \frac{3 \pi c^2}{2 \omega_0^3} \frac{\Gamma}{\Delta} I,
\end{equation}
for the ground and excited state (upper and lower sign,
respectively).
Equation \eqref{eq:light_shift_2lv_atoms} shows that the energy shift of the ground state (the solution with the sign $+$) exactly correspond to the dipole potential \eqref{eq:dipole_potential}. This allows the reinterpretation of the semiclassical result obtain for the external degrees of freedom (position and momentum) in terms of the dynamic of the internal degrees of freedom. In presence of the laser field, the energy levels are shifted by a quantity proportional to the field intensity with a sign that depends on the sign of the detuning. Assuming a low saturation regime, corresponding to $s\ll1$, where the saturation parameter $s$ is given by
\begin{equation}
    s = \frac{\Omega^2/2}{\Delta^2 +(\Gamma^2/4)},
\end{equation}
and $\Omega$ is the Rabi frequency, defined as
\begin{equation}
    \Omega = \frac{|\mathcal{E}_0|}{\hbar}\braket{e |\,\hat{d}\,| g},
\end{equation}
the atomic population can be considered almost entirely confined to the ground state. Consequently, the light-shifted ground state becomes the effective potential that governs the motion of the atoms, driving them toward regions of minima of $\Delta E$ (Figure \ref{img:light_shift}). 

\begin{figure}
    \centering
    \includegraphics[width=0.7\linewidth]{Cap_1_figures/light-shif.png}
    \caption{Left: light-shifted energy levels for a two-level atom for a red detuning ($\Delta < 0$). Right: a spatially inhomogeneous field produces a potential well that traps the atoms. Credits \cite{grimm}.}
    \label{img:light_shift}
\end{figure}

For a multi-level atom Equation~\eqref{eq:energy_perturbative_correction} contains many terms similar to equation \eqref{eq:light_shift_2lv_atoms}, accounting for all the allowed transitions from any ground state sublevel $\ket{g_i}$ to any excited state $\ket{e_j}$. Each of these terms depends on the dipole matrix element $d_{ij}=\braket{g_i|\,\hat{d}\,|e_j}$. Applying the Wigner–Eckart theorem \cite{sakurai2017modern}, this matrix element can be factored as:
\begin{equation}
    d_{ij}=c_{ij}||d||
\end{equation}
where $||d||$ is the reduced matrix element given by \eqref{eq:quantum_gamma} and $c_{ij}$ are real Clebsch-Gordan coefficients describing the selection rules and coupling strengths of the transitions. They depend on the light polarization and the angular momentum quantum numbers of the involved hyperfine states.

Therefore, the  energy shift \eqref{eq:energy_perturbative_correction} for the ground state $\ket{g_i}$ is
\begin{equation}
    \label{eq:light_shift_multi_lv_atoms}
    \Delta E_i = \frac{3 \pi c^2 \Gamma}{2 \omega_0^3} I \times \sum_j \frac{c_{ij}^2}{\Delta_{ij}}
\end{equation} 
and according to the interpretation given from the two-level atoms, the dipole potential at low saturation is
\begin{equation}
    U_\text{dip} = \Delta E_i.
\end{equation}
Thus, the dipole potential in a multi-level atom arises from the contribution of all the allowed transitions, each weighted by the ratio between its line strengths $c_{ij}^2$ and its detuning $\Delta_{ij}$.

\subsubsection{Fine and Hyperfine structures contribution for Alkali atoms}
Many laser cooling experiments utilize alkali metal atoms. This preference arises because their principal transition frequencies fall within the visible and infrared (VIS-IR) range making them easily accessible with standard laser technologies. Moreover, their relatively simple atomic structures is particularly convenient, allowing easy identification of closed transitions necessary for continuous cooling cycles, while minimizing the number of repumper lasers needed to address the population leaking to dark states.

When considering the fine structure arising from the spin-orbit interaction, the electronic ground state of alkali atoms is $^2S_{1/2}$. From this state, the electric dipole selection rules ($\Delta L = \pm 1$, $\Delta S = 0$) allow for two primary transitions, known as $D-$lines:
\begin{itemize}
    \item $D_1$: $^2S_{1/2} \rightarrow \, ^2P_{1/2}$;
    \item $D_2$: $^2S_{1/2} \rightarrow \, ^2P_{3/2}$.
\end{itemize}
Beyond the fine structure, the energy levels undergo further splitting due to the coupling between the nuclear spin and the total electronic angular momentum. The resulting configuration is called hyperfine structure. Denoting the fine structure splitting as $\hbar\Delta_{FS}$, the ground state hyperfine splitting as $\hbar\Delta_{HFS}$, and the excited state hyperfine splitting as $\hbar\Delta'_{HFS}$, these characteristic energy scales obey $\Delta_{FS}\gg\Delta_{HFS}\gg\Delta'_{HFS}$. The complete electronic energy level scheme for an alkali atom with nuclear spin $I=3/2$ is shown in figure \ref{fig:alkali_electronic_energy_levels}.
\begin{figure}
    \centering
    \includegraphics[width=0.55\linewidth]{Cap_1_figures/alkali_energy_levels.png}
    \caption{Energy level scheme of an alkali atom with $I=3/2$ such as $^{87}Rb$. Credits \cite{grimm}.}
    \label{fig:alkali_electronic_energy_levels}
\end{figure}

Assuming that all optical detunings stay large compared with the excited-state hyperfine splitting $\Delta'_{HFS}$, equation \eqref{eq:light_shift_multi_lv_atoms} can be calculated for an alkali atom in the ground state of total angular momentum $F$ and magnetic quantum number $m_F$, obtaining \cite{grimm}:
\begin{equation}
    \label{eq:alkali_dipole_potential_general_case}
    U_{\text{dip}}(\mathbf{r}) = \frac{\pi c^2 \Gamma}{2 \omega_0^3} \left( \frac{2 + \mathcal{P} g_F m_F}{\Delta_{2,F}} + \frac{1 - \mathcal{P} g_F m_F}{\Delta_{1,F}} \right) I(\mathbf{r}),
\end{equation}
where the two term in the parentheses represent respectively the contribution of the $D_2$ and $D_1$ line. Here, $g_F$ is the Landé factor, $\mathcal{P}$ characterizes the laser polarization ($\mathcal{P}=0,\pm1$ for linear polarization and circular $\sigma^\pm$ respectively), and $\Delta_{2,F}$ and $\Delta_{1,F}$ denote the detunings of the laser field relative to the transitions from the $^2S_{1/2}, F$ ground state to the center of respectively the $^2P_{3/2}$ and $^2P_{1/2}$ hyperfine excited state.

\vspace{0.5\baselineskip}
If the detunings are large with respect to the fine-structure splitting ($|\Delta_{1,F}|, |\Delta_{2_F}| \gg \Delta'_{FS}$), introducing the detuning $\Delta$ with respect to the center of the $D-$line doublet, expanding equation \eqref{eq:alkali_dipole_potential_general_case} at first order in the small parameter $\Delta'_{FS}/\Delta$, one gets \cite{grimm}:
\begin{equation}
    \label{eq:alkali_dipole_potential_large_detuning}
    U_{\text{dip}}(\mathbf{r}) = \frac{3 \pi c^2}{2 \omega_0^3} \frac{\Gamma}{\Delta} \left( 1 + \frac{1}{3} \mathcal{P} g_F m_F \frac{\Delta'_{\text{FS}}}{\Delta} \right) I(\mathbf{r}).
\end{equation}
The zeroth order term corresponds exactly to the dipole potential of a simple two-level atom (see Eq. \eqref{eq:dipole_potential}). The first-order term introduces a small, spin-dependent correction depending on the laser polarization $\mathcal{P}$ and the magnetic quantum number $m_F$.

Physically, equation \eqref{eq:alkali_dipole_potential_large_detuning} describes that, when the fine structure is completely unresolved due to the large laser detuning, the atom effectively acts as a generic two-level system undergoing an $s \rightarrow p$ transition.
Consequently, the resulting dipole potential is dominated by the scalar term, matching the two-level model previously discussed.

Since in the unresolved fine structure case, the polarization-dependent term is heavily suppressed by the small parameter $\Delta'_{FS}/\Delta$ (and vanishes entirely for linearly polarized light), the atom effectively acquires only the scalar light shift of the unperturbed $s$ state. Consequently, all hyperfine and magnetic sublevels experience a uniform trapping potential that confines the atom while preserving its internal spin degeneracy.

In a more general case, when the fine structure is resolved, but the hyperfine structure is unresolved ($|\Delta_{1,F}|,|\Delta_{2,F}|\gg \Delta_{HFS} \gg \Delta'_{HFS}$), only linearly polarized light guarantees that the dipole potential becomes completely independent of $m_F$ and $F$ \cite{grimm}. Therefore, linearly polarized light is usually preferred for dipole trap implementation.


%* ##########################################################
\subsection{Heating mechanisms}
%* ##########################################################

Typical trap depths for optical dipole traps are usually below 1 mK. Therefore, capturing a substantial number of atoms requires a preliminary cooling stage. For instance, the experiment discussed in this thesis employs a magneto-optical trap (MOT) followed by an optical molasses phase.

Given these shallow confining potentials, even minor heating rates can largely limit the trap lifetime. It is therefore important to outline the mechanisms responsible for the parasitic heating induced by the trapping light.

\subsubsection{Heating source}
In general, when atoms interact with light and absorb photons, when the subsequent spontaneous emission event occur, the randomness of the emitted photons' direction causes a diffusion in the phase space that corresponds to a heating processes. This discrete momentum transfer is the fundamental phenomenon that establishes lower limits for laser cooling mechanisms \cite{CohenLesHouches, metcalf1999}.

At large detuning the scattering process follows a Poisson statistic. In this regime, the absorption of a photon from the dipole beam introduces a heating energy equivalent to the recoil energy $E_\text{rec} = k_B T_\text{rec}/2$ per scattering event, where $T_\text{rec}$ is the recoil temperature. Since the absorptions occur strictly along the propagation axis the momentum diffusion occurs only along this axis, making the associated heating effect anisotropic. On the other hand, the subsequent spontaneous emission also increase the thermal energy by one recoil energy $E_\text{rec}$ per event, however the emission directions are random, thus the effect is spread over all the three spatial direction.

Therefore, in a single scattering cycle, the atom gains $E_\text{rec}$ of thermal energy from the absorption and then $E_\text{rec}$ from the spontaneous emission, for a total of $2E_\text{rec}$, in a time $\bar{\Gamma}_\text{sc}^{-1}$. It leads to the total heating power
\begin{equation}
    \label{eq:heating_power}
    P_{\text{heat}} = 2 E_{\text{rec}} \bar{\Gamma}_{\text{sc}} = k_B T_{\text{rec}} \bar{\Gamma}_{\text{sc}},
\end{equation}
the heating depends on the mean scattering rate $\bar{\Gamma}_\text{sc}$ and on the recoil $E_\text{rec}$ which depends on the laser frequency.

Other fundamental source of heating, such as the redistribution of photons between different travelling wave, should be quenched by the far-detuned working frequency \cite{grimm}.

In addition to the fundamental heating mechanism, it must be mentioned also the presence of technical heating caused by fluctuations in the laser power and in its pointing direction. In particular power fluctuations at twice the characteristic trap frequencies are relevant, as they can parametrically drive the oscillatory atomic motion \cite{Savard:97}.

\subsubsection{Heating rate}
Assuming that the trapped atoms are at thermal equilibrium, the mean energy can be expressed through the energy equipartition theorem. Assuming the trap as a three-dimensional harmonic potential well, there is a contribution of three quadratic degrees of freedom from the kinetic energy, and other three from the harmonic potential. According to the energy partition theorem, the mean energy is
\begin{equation}
    \label{eq:equipartition_th}
    \bar{E}= 3 k_B T.
\end{equation}
Using \eqref{eq:heating_power} in \eqref{eq:equipartition_th}, one can obtain an equation that describes the temperature change with time:
\begin{equation}
    \label{eq:temperature_changing_rate}
    \dot{T} = \frac{1}{3} T_\text{rec} \bar{\Gamma}_\text{sc}.
\end{equation}

Expressing the dipole potential as
\begin{equation}
    \label{eq:mean_scattering_rate_td}
    \bar{U}_\text{dip} = U_0 + \bar{E}_\text{pot} = U_0 + \frac{3}{2}k_B T,
\end{equation}
where $U_0$ is the potential energy minimum.

Using equation \eqref{eq:dipole_scattering}, $\bar{U}_\text{dip}$ can be expressed as:
\begin{equation}
    \bar{U}_\text{dip} = \frac{\hbar \Delta \Gamma_\text{sc}}{\Gamma}.
\end{equation}
Substituting this into Equation \eqref{eq:mean_scattering_rate_td} yields the mean scattering rate:
\begin{equation}
    \label{eq:scattering_rate_Tequilibrium}
    \bar{\Gamma}_\text{sc} = \frac{\Gamma}{\hbar\Delta} (U_0 + \frac{3}{2}k_B T).
\end{equation}
It is important to notice that, for a blue-detuned trap, $U_0 = 0$ and the trap depth $\hat{U}$ is given by the maximum value of the potential surrounding the minimum, by contrast for a red-detuned trap $\hat{U}=|U_0|$.
Finally, using equation \eqref{eq:scattering_rate_Tequilibrium} in \eqref{eq:temperature_changing_rate}, with a depth $\hat{U} \gg k_BT$, one gets the following temperature changing rate for a red and a blue trap:
\begin{subequations}
    \label{eq:t_rate_blue_red_comparison}
    \begin{align}
        \dot{T}_{\text{red}} &= \frac{1}{3} T_{\text{rec}} \frac{\Gamma}{\hbar |\Delta|} \hat{U}, \\[2mm]
        \dot{T}_{\text{blue}} &= \frac{1}{2} T_{\text{rec}} \frac{\Gamma}{\hbar \Delta} k_B T.
    \end{align}
\end{subequations}
Equations \eqref{eq:t_rate_blue_red_comparison} describe the increasing in temperature of the atomic cloud if subjected to the dipole trap only. It highlights another blue trap limitation: while a red trap has a linear heating rate (as long as $\hat{U}\gg k_BT$), the blue trap shows an exponential heating law.